\section{Automatic Gain Control Circuit}

\subsection{Basic blocks}
\begin{figure}[H]
    \centering
    % Left side: explanation + equations
    \begin{minipage}[c]{0.55 \linewidth}
    In audio processing, we want to modify and control the amplitude of a signal, enhancing or cutting the dynamics of low-amplitude (fader, expander) and/or high-amplitude (compressor, limiter) audio signals.  
The idea is to use a \textbf{Voltage Controlled Amplifier (VCA)} which adjusts the amplification and gain based on a reference voltage $V_c$.
\[
v_{out} = G(V_c) \cdot v_{in}
\]
    \end{minipage}\hfill
    \begin{minipage}[c]{0.45\linewidth}
        \centering
        \includegraphics[width=\linewidth]{VCA.png}
    \caption{VCA block and transfer function.}
    \end{minipage}\hfill
\end{figure}

To implement this, one building block is the \textbf{voltage multiplier}:
\[
v_{out} = v_a \cdot v_b
\]
\begin{figure}[H]
    \centering
    \includegraphics[width=0.7\linewidth]{Voltage_Multiplier.png}
    \caption{Voltage Multiplier block representation.}
\end{figure}
To compute the signal’s average, we define the \textbf{integrator}:
\[
\overline{v} = \frac{1}{T} \int_T v(\tau) \, d\tau
\]
In OPAMP form:
\[
v_{out}(t) = -\frac{1}{RC} \int v_{in}(t) \, dt
\]
\begin{figure}[H]
    \centering
    \includegraphics[width=0.7\linewidth]{integrator.png}
    \caption{Integrator block representation.}
\end{figure}
To compute the \textbf{signal RMS}, we also define the \textbf{square root} operator.
\begin{figure}[H]
    \centering
    % Left side: explanation + equations
    \begin{minipage}[c]{0.55 \linewidth}
    First, we define:
    \[
    v_2 = v_{out}^2
    \]
    From the ideal OPAMP laws:
    \[
    i_1 + i_2 = i^- = 0
    \]
    with
    \[
    i_1 = \frac{v_{in}}{R_1}, 
    \qquad i_2 = \frac{v_2}{R_2} = \frac{v_{out}^2}{R_2}
    \]
    Thus:
    \[
    \frac{v_{in}}{R_1} + \frac{v_{out}^2}{R_2} = 0 
    \quad \Rightarrow \quad
    v_{out} = \sqrt{-\frac{R_2}{R_1} v_{in}}
    \]
    \end{minipage}\hfill
    \begin{minipage}[c]{0.45\linewidth}
        \centering
        \includegraphics[width=\linewidth]{square_root.png}
    \caption{Circuit for the square root operator.}
    \end{minipage}\hfill
\end{figure}

\textbf{Note:} The formula gives real results only if $v_{in} < 0$, hence we use an inverting buffer (as in the integrator).  

The complete signal path to compute the RMS is:
\begin{figure}[H]
    \centering
    \includegraphics[width=0.7\linewidth]{RMS_Pipeline.png}
    \caption{RMS computation block diagram.}
    \label{fig:placeholder}
\end{figure}

For a sinusoidal input $v_{in}$, the RMS measurement is:
\[
v_{RMS} = \frac{v_{peak}}{\sqrt{2}}
\]
\begin{figure}[H]
    \centering
    \includegraphics[width=0.7\linewidth]{RMS_sine.png}
    \caption{Sine signal RMS computation illustration.}
\end{figure}

\subsection{Compressor and Limiter}

The Automatic Gain Control requires a voltage reference, which depends directly on the input signal itself.  

\begin{figure}[H]
    \centering
    \includegraphics[width=0.8\linewidth]{AGC_Circuit.png}
    \caption{Illustration of the voltage threshold inclusion in the circuit.}
\end{figure}

The \textbf{threshold condition} is:
\[
v_{in} > V_{threshold} \quad \Rightarrow \quad \text{signal is attenuated (compression/limiting).}
\]

\begin{figure}[H]
    \centering
    % Left side: explanation + equations
    \begin{minipage}[c]{0.35 \linewidth}
    \centering
    \includegraphics[width=\linewidth]{linear_scale.png}
    \caption{Linear scaling of the compressor/limiter.}
    \end{minipage}\hfill
    \begin{minipage}[c]{0.6\linewidth}
        \centering
        \includegraphics[width=\linewidth]{Limiter.png}
    \caption{Attack and release times highlighted.}
    \end{minipage}\hfill
\end{figure}



\begin{figure}[H]
    \centering
    % Left side: explanation + equations
    \begin{minipage}[c]{0.6 \linewidth}
    Gain control is not immediate. We define:
\begin{itemize}
    \item \textbf{Attack time}: the time needed for the system to reduce the gain when the input amplitude suddenly increases.
    \item \textbf{Release time}: the time needed for the system to restore the gain after the input amplitude decreases.
\end{itemize}

These times determine how the gain changes during transitions between low and high amplitudes.  

The RMS-based description of a compressor/limiter can be expressed as: $ v_{out,RMS} \quad \text{versus} \quad v_{in,RMS}$
    \end{minipage}\hfill
    \begin{minipage}[c]{0.38\linewidth}
        \centering
        \includegraphics[width=\linewidth]{limiter_comparison.png}
    \caption{Comparison in between free compression and compressed signal.}
    \end{minipage}\hfill
\end{figure}

The degree of compression is described by the \textbf{Compression Ratio} $1:k$. This means:
\[
\text{Input increase of } k \, \text{dBu} \quad \Rightarrow \quad \text{Output increase of } 1 \, \text{dBu}
\]

\begin{figure}[H]
    \centering
    % Left side: explanation + equations
    \begin{minipage}[c]{0.5 \linewidth}
    We distinguish 2 types of transitions:
\begin{itemize}
    \item \textbf{Hard knee} (red curve): the compression starts abruptly at the threshold.  
    \item \textbf{Soft knee} (blue curve): the compression starts gradually around the threshold.  
\end{itemize}
    \end{minipage}\hfill
    \begin{minipage}[c]{0.5\linewidth}
        \centering
        \includegraphics[width=\linewidth]{Hard_Knee_Soft_Knee.png}
    \caption{Hard knee (left) and Soft knee (right).}
    \end{minipage}\hfill
\end{figure}

\textbf{Limiter vs Clipper:}
    \begin{itemize}
        \item A \textbf{limiter}, using gain control, prevents the signal from exceeding a threshold but does not introduce distortion.  
        \item A \textbf{clipper}, instead, abruptly cuts the signal above a threshold, introducing distortion.  
    \end{itemize}
\begin{figure}[H]
    \centering
    \includegraphics[width=0.3\linewidth]{Limiter_vs_Clipper.png}
    \caption{Limiter and Clipper comparison.}
\end{figure}

Based on the temporal evolution of the input, the gain control function becomes:
\[
G(t) = f\big(v_{in}(t), \, V_{threshold}, \, \text{attack time}, \, \text{release time}, \, k \big)
\]
where $k$ is the compression ratio.
\begin{figure}[H]
    \centering
    \includegraphics[width=0.4\linewidth]{Gain_Control.png}
    \caption{Gain control as a function of time.}
    \label{fig:placeholder}
\end{figure}

\subsection{Gate and Expander}

\begin{figure}[H]
    \centering
    % Left side: explanation + equations
    \begin{minipage}[c]{0.65 \linewidth}
    We now focus on controlling low amplitudes.  
\begin{itemize}
    \item The \textbf{Gate} is used to \textit{cut signals} that are lower than a threshold.  
    \item The \textbf{Expander} is used to \textit{change the gain} (and therefore the compression ratio) for low-level signals.  
\end{itemize}

The structure is similar to the Compressor/Limiter, but this time we switch the lower OPAMP inputs: RMS on pin "-" and $V_{threshold}$ on pin "+".

We have: 
\[
v_{out,RMS} = f(v_{in,RMS}, \, V_{threshold}, \, k_{exp})
\]
where $k_{exp}$ is the expansion ratio.  
    \end{minipage}\hfill
    \begin{minipage}[c]{0.3\linewidth}
        \centering
        \includegraphics[width=\linewidth]{Gate_Expander_Circuit.png}
    \caption{Circuit of the Gate/Expander system.}
    \end{minipage}\hfill
\end{figure}

\textbf{Gate:}
If $v_{in,RMS} < V_{threshold}$, the output is attenuated strongly (or cut).  
This suppresses background noise or unwanted weak signals.

\textbf{Expander:}
The expander reduces, but does not cut, low signals.  
This results in an \textbf{enhanced dynamic range}, since the difference between low and high levels is increased.

\begin{figure}[H]
    \centering
    \includegraphics[width=0.75\linewidth]{Gate_vs_Expander.png}
    \caption{Gate (left) and Expander (right).}
\end{figure}


\textbf{Examples of Dynamic Control:}
\begin{itemize}
    \item A \textbf{Downward Expander} reduces the volume of signals below the threshold.
    \item An \textbf{Upward Expander} increases the level of signals above the threshold, expanding dynamics in the opposite way.
\end{itemize}

\begin{figure}[H]
    \centering
    \includegraphics[width=0.5\linewidth]{Expander_Dynamic_Control.png}
    \caption{Dynamic control modification of the expander.}
\end{figure}

\textbf{Related Techniques:}
\begin{itemize}
    \item \textbf{De-esser:} Works like a compressor, but limits compression to a specific frequency band (e.g., sibilant sounds in vocals).  
    \item \textbf{Side-chain compression:} The threshold is controlled by an \textit{external signal}.  
    Example: lowering the music when a voice is detected (common in broadcasting).  
\end{itemize}

